\newpage
\noindent
Since, eliminating $t$, $y_2\propto y_1^{\lambda_2/\lambda_1}
=\dfrac{1}{y_1^{|\lambda_2/\lambda_1|}}$.
So, orbits are \underline{hyperbolas}!

\bigskip
\noindent
\textcircled{2}\quad\underline{Node}:
$\underline{(\operatorname{Tr}(\hat A))^2>4\operatorname{Det}(\hat A)>0}$,
$\Rightarrow\lambda_1,\lambda_2$ are \underline{real} and either both positive
$(\operatorname{Tr}(\hat A)>0)$ or negative $(\operatorname{Tr}(\hat A)<0)$.

\bigskip
\noindent
\begin{minipage}[t]{0.50\textwidth}
\vspace{0pt}
Same canonical form
\[
\hat A=\begin{pmatrix}\lambda_1&0\\0&\lambda_2\end{pmatrix},
\quad\Rightarrow\quad\underline{y_2\propto y_1^{\lambda_2/\lambda_1}}.
\]
Get different \underline{patterns} depending on
$\underline{\lambda_2/\lambda_1<1\text{ or }>1}$.

\bigskip
\underline{Sign of $\lambda_1,\lambda_2$ sets the}
\underline{direction of the flow}.
\end{minipage}\hfill
\begin{minipage}[t]{0.46\textwidth}
\vspace{0pt}
\includegraphics[width=\linewidth]{lecture5-unstable-nodes.png}
\end{minipage}

\begin{center}
\includegraphics[width=0.78\textwidth]{lecture5-stable-nodes.png}
\end{center}

\newpage
\noindent
\textcircled{3}\quad\underline{Focus (spiral)}:
$\underline{(\operatorname{Tr}(\hat A))^2<4\operatorname{Det}(\hat A)}$
and $\operatorname{det}(\hat A)>0$ so that roots are complex,
$\underline{\lambda=\rho\pm i\omega}$, $\omega\neq0$.

\bigskip
\noindent
\begin{minipage}[c]{0.49\textwidth}
Canonical form
\[
\hat A=\begin{pmatrix}\rho&-\omega\\\omega&\rho\end{pmatrix}.
\]
In polars have
\[
\dot r=\rho r,\quad\dot\theta=\omega,\quad\Rightarrow
\]
\[
\underline{r=r_0e^{\rho t},\quad\theta=\theta_0+\omega t}.
\]
\end{minipage}\hfill
\begin{minipage}[c]{0.49\textwidth}
\includegraphics[width=\linewidth]{lecture5-focus.png}
\end{minipage}

\bigskip
\noindent Other cases are \underline{degenerate}.

\bigskip
\noindent
\textcircled{4}\quad
$\underline{(\operatorname{Tr}(\hat A))^2=4\operatorname{Det}(\hat A)>0}$

\medskip
\noindent
A boundary between \underline{focus \& node}, double eigenvalue
$\underline{\lambda=\frac12\operatorname{Tr}(\hat A)}$. Two cases:

\bigskip
\noindent
\begin{minipage}[c]{0.49\textwidth}
\underline{Stellar node (star)}
\[
\text{if }\hat A=\begin{pmatrix}\lambda&0\\0&\lambda\end{pmatrix}
\text{ is canonical}
\]
\end{minipage}\hfill
\begin{minipage}[c]{0.49\textwidth}
\includegraphics[width=\linewidth]{lecture5-stellar-nodes.png}
\end{minipage}

\newpage
\noindent
\begin{minipage}[c]{0.48\textwidth}
\underline{Improper node}
\[
\text{if }\hat A=\begin{pmatrix}\lambda&1\\0&\lambda\end{pmatrix}
\text{ is canonical}
\]
\centerline{(defective)}
\[
\text{Here}\quad y=y_0e^{\lambda t},
\]
\[
x=x_0e^{\lambda t}+y_0te^{\lambda t},
\qquad\frac{y}{x}=\frac{y_0}{x_0+y_0t}.
\]
\end{minipage}\hfill
\begin{minipage}[c]{0.50\textwidth}
\includegraphics[width=\linewidth]{lecture5-improper-node.png}
\end{minipage}

\bigskip
\noindent
\textcircled{5}\quad
$\underline{\operatorname{Det}(\hat A)>0,\quad\operatorname{Tr}(\hat A)=0}$:
\underline{Centre}

\medskip
\noindent
\begin{minipage}[c]{0.59\textwidth}
\[
\text{Here}\quad\lambda=\pm i\omega,\qquad
\hat A=\begin{pmatrix}0&-\omega\\\omega&0\end{pmatrix}.
\]
\[
\underline{\dot r=0},\quad\underline{\dot\theta=\omega},
\quad\Rightarrow\quad\text{orbits are closed circles}.
\]
\end{minipage}\hfill
\begin{minipage}[c]{0.38\textwidth}
\includegraphics[width=\linewidth]{lecture5-centre.png}
\end{minipage}

\bigskip
\noindent
\textcircled{6}\quad
$\underline{\operatorname{Det}(\hat A)=0,\quad\operatorname{Tr}(\hat A)\neq0}$
-- between \underline{saddle \& node}

\medskip
\noindent
\begin{minipage}[c]{0.48\textwidth}
\[
\lambda_1=0,\quad\lambda_2\neq0,\qquad
\hat A=\begin{pmatrix}0&0\\0&\lambda\end{pmatrix}
\]
\[
y_1=\operatorname{const},\qquad y_2=y_{2,0}e^{\lambda t}
\]
\hfill for $\lambda<0$
\end{minipage}\hfill
\begin{minipage}[c]{0.49\textwidth}
\includegraphics[width=\linewidth]{lecture5-line-fixed-points.png}
\end{minipage}

\newpage
\noindent
\textcircled{7}\quad
$\underline{\operatorname{Det}(\hat A)=0,\quad\operatorname{Tr}(\hat A)=0}$,
$\lambda_{1,2}=0$

\medskip
\noindent
if $\hat A=0$, $\Rightarrow y_1=\operatorname{const},\ y_2=\operatorname{const}$,
$\Rightarrow$ plane of fixed points!

\medskip
\noindent
\begin{minipage}[c]{0.45\textwidth}
if $\hat A=\begin{pmatrix}0&1\\0&0\end{pmatrix}$, this is a case of shear,
\[
\underline{\dot y_1=y_2,\quad\dot y_2=0}.
\]
\end{minipage}\hfill
\begin{minipage}[c]{0.53\textwidth}
\includegraphics[width=\linewidth]{lecture5-shear.png}
\end{minipage}

\medskip
\noindent This zoo can be summarized in a diagram
\begin{center}
\includegraphics[width=0.96\textwidth]{lecture5-classification.png}
\end{center}

\noindent Canonical forms are in special basis
\[
\begin{array}{cll}
(\ast)&\begin{pmatrix}\lambda_1&0\\0&\lambda_2\end{pmatrix}
&\text{-- eigenvectors}\\[12pt]
(\ast)&\begin{pmatrix}\lambda&1\\0&\lambda\end{pmatrix}
&\text{-- generalized eigenvectors}\\[12pt]
(\ast)&\begin{pmatrix}\rho&-\omega\\\omega&\rho\end{pmatrix}
&\text{-- Re \& Im of eigenvectors}
\end{array}
\]

\newpage
\noindent\underline{Notes}

\medskip
\noindent
$(\ast)$ No need to change basis in practice, just work with
$\operatorname{Tr}(\hat A)$, $\operatorname{Det}(\hat A)$ to classify

\medskip
\noindent
$(\ast)$ In original (non-canonical) coordinates shapes can be distorted
\begin{center}
\includegraphics[width=\textwidth]{lecture5-distorted.png}
\end{center}

\noindent
$(\ast)$ Useful to know eigenvectors for saddles/nodes; for focus/centre
cases ratios depend also on $\lambda_{1,2}$, easier to just plot several values

\medskip
\noindent
$(\ast)$ Centres are rare but generic in Hamiltonian systems:
\[
\dot{\bar x}=\begin{pmatrix}\dfrac{\partial H}{\partial y}\\[6pt]
-\dfrac{\partial H}{\partial x}\end{pmatrix}
\quad\text{for some }H(x,y),\quad\bar x=(x,y).
\]
\[
\text{Indeed,}\quad
\hat A=\begin{pmatrix}H_{xy}&H_{yy}\\-H_{xx}&-H_{xy}\end{pmatrix},
\quad\Rightarrow\quad\underline{\operatorname{Tr}(\hat A)=0},\quad\Rightarrow
\]
\[
\text{either}\quad
\underset{\lambda=\pm i\omega}{\underline{\text{centres}}}
\quad\text{or}\quad
\underset{\lambda_1\lambda_2<0}{\underline{\text{saddles}}}.
\]
\[
\text{Also,}\quad
\underline{\frac{dH}{dt}}
=(\dot{\bar x}\cdot\nabla)H
=\dot xH_x+\dot yH_y=H_yH_x-H_xH_y=0.
\]
\noindent
Thus, $H(x,y)=\operatorname{const}$ on trajectories, orbits follow contours
of $H(x,y)$ in \underline{autonomous Hamiltonian} systems.

\begin{center}
\textbf{End of Lecture 5}
\end{center}
